\documentclass[11pt]{article}

\usepackage[margin=1in]{geometry}
\usepackage{amsmath,amssymb,amsthm,mathtools,microtype}
\usepackage[T1]{fontenc}
\usepackage[hidelinks]{hyperref}

\newtheorem{theorem}{Theorem}[section]
\newtheorem{remark}[theorem]{Remark}

\newcommand{\C}{\mathbb C}
\newcommand{\N}{\mathbb N}
\newcommand{\E}{\mathbb E}
\newcommand{\Supp}{\operatorname{Supp}}

\title{Sparse Minimality of Gaussian-Moments Counterexamples\\
in Dimension Three}
\author{Roy van Rijn\\{\small Independent researcher}}
\date{17 November 2026}
\hypersetup{
  pdftitle={Sparse Minimality of Gaussian-Moments Counterexamples in Dimension Three},
  pdfauthor={Roy van Rijn},
  pdfsubject={Sparse minimal counterexamples to the Gaussian Moments Conjecture in dimension three},
  pdfkeywords={Gaussian Moments Conjecture, Gaussian moments,
    sparse polynomials, Groebner bases, computer-assisted proof}
}

\begin{document}
\maketitle

\begin{abstract}
We prove that a polynomial in three real Gaussian variables with total
degree at most four and at most four nonzero monomials in a fixed Gaussian
Witt frame satisfies the mixed-moment conclusion whenever all its positive
pure moments vanish.  The proof combines an exhaustive exact calculation
over \(59{,}535\) supports with all-order arguments for three surviving
coefficient families: isotropic linear forms, polynomial Gaussian shears,
and orthogonal products with an isotropic factor.  Moment ideals through
order twenty-four exclude every other active support.  Together with
Long's five-term quartic counterexample, the result shows that five
Witt-coordinate monomials are necessary and sufficient for a failure of
degree at most four.  The degree bound and the chosen notion of support
are essential restrictions of the result.
\end{abstract}

\section{Introduction}

Let \(X_1,\ldots,X_d\) be independent standard real Gaussian random
variables.  The Gaussian Moments Conjecture, introduced by Derksen, van den
Essen, and Zhao~\cite{DVZ}, asserts that if
\[
  P\in\C[X_1,\ldots,X_d],
  \qquad \E(P^m)=0\quad(m\geq1),
\]
then, for every \(Q\in\C[X_1,\ldots,X_d]\),
\[
  \E(QP^m)=0
  \qquad\text{for all sufficiently large }m.
\]
Long constructed explicit counterexamples in every dimension
\(d\geq3\)~\cite[Theorem~5.1 and Corollary~5.2]{Long}.  His three-variable witness, in circular
coordinates \(W,Z\) and an independent real Gaussian \(T\), is
\begin{equation}\label{eq:long}
 P_{\mathrm L}
 =W+WZ-T^2-\frac32ZT^2-\frac12Z^2T^2.
\end{equation}
It has degree four and five monomials, all its pure moments vanish, and
\(\E(ZP_{\mathrm L}^m)=m!\) for every \(m\geq1\).  We ask whether a failure of degree at most four can use fewer monomials.
Support size throughout refers to a fixed Gaussian Witt frame; it is not
an invariant of arbitrary polynomial coordinate changes.

\begin{samepage}
\begin{theorem}[Bounded sparse exclusion]\label{thm:sparse}
Let \(X,Y,T\) be independent standard real Gaussians and set
\[
 Z=(X+iY)/\sqrt2,\qquad W=(X-iY)/\sqrt2.
\]
These Gaussian Witt coordinates satisfy
\[
 \E(W^iZ^jT^k)=
 \begin{cases}
   i!(k-1)!!,&i=j\text{ and }k\text{ is even},\\
   0,&\text{otherwise},
 \end{cases}
\qquad (-1)!!=1.
\]
Suppose that \(P\in\C[W,Z,T]\) has total degree at most four and at most four
nonzero monomials.  If \(\E(P^m)=0\) for every \(m\geq1\), then, for every
\(Q\in\C[W,Z,T]\),
\[
 \E(QP^m)=0
 \qquad\text{for all sufficiently large }m.
\]
Consequently no GMC counterexample has both degree at most four and support
size at most four in a Gaussian Witt frame.
\end{theorem}
\end{samepage}

Thus five terms are necessary for a counterexample of degree at most
four, and \eqref{eq:long} attains this term bound.  This statement does not
prove global degree minimality or Pareto minimality of the pair \((4,5)\):
a five-term cubic would lie outside the excluded rectangle.  A four-term
counterexample of degree greater than four is also outside the theorem.

\paragraph{Why finite equations can settle this bounded problem.}
Gaussian expectation can only see products whose \(W\)- and \(Z\)-degrees
balance and whose \(T\)-degree is even.  For example, if
\(P=aW+bZ\), then \(\E(P^2)=2ab\); when both coefficients are nonzero,
vanishing already fails at the second moment.  More complicated supports
can have several balanced products at the same order, allowing their
coefficient contributions to cancel.  The finite calculation must
therefore solve coefficient equations, not merely look for balanced
products.

There are only \(35\) monomials of degree at most four.  For each small
support, exact elimination either rules out such cancellation or produces
one of three explicit families.  The difficult logical transition is from
these finite equations to all sufficiently large mixed moments.  It is
made separately: each surviving family is carried to a polynomial with
one-sided isotropic weight by a Gaussian-preserving change of variables.
A fixed multiplier has bounded weight and cannot balance arbitrarily
large powers.  No assertion that twenty-four moments determine all
moments is needed.

\section{Moment ideals}

The finite part of the proof is an exhaustive exact coefficient-scheme
calculation; the all-order interpretation of its survivors is conceptual.
Let
\[
 \mathcal M_4
 =\{W^iZ^jT^k:i,j,k\geq0,\ i+j+k\leq4\}.
\]
This set has \(35\) elements.  Fix a support \(S\subset\mathcal M_4\), write
\[
 P_S=\sum_{s\in S}c_ss,
\]
and work on the coefficient torus \(\prod_{s\in S}c_s\neq0\).  For
\(s=(i_s,j_s,k_s)\) and a multiplicity vector
\(\nu=(\nu_s)_{s\in S}\in\mathbb N^S\), put
\[
 |\nu|=\sum_s\nu_s,\quad
 I(\nu)=\sum_s\nu_si_s,\quad
 J(\nu)=\sum_s\nu_sj_s,\quad
K(\nu)=\sum_s\nu_sk_s.
\]
Before expanding any Wick coefficients, associate to \(S\) its contraction
semigroup
\[
 \Gamma_S=
 \left\{\nu\in\mathbb N^S:
   \sum_s\nu_s(i_s-j_s)=0,\quad
   \sum_s\nu_sk_s\equiv0\pmod2
 \right\}.
\]
Its Hilbert basis \(H_S\) consists of the primitive contractions.  This
basis is computed exactly in two stages.  First ignore parity and
enumerate the indecomposable nonnegative solutions of
\(\sum_s q_s\nu_s=0\), where \(q_s=i_s-j_s\).  The following elementary
one-row bound makes this enumeration finite:
\begin{equation}\label{eq:charge-bound}
 |\nu|\leq 1+\sum_{s\in S}|q_s|.
\end{equation}
Indeed, a zero-charge monomial contributes only its unit vector as an
indecomposable solution.  Otherwise an indecomposable solution uses
positive and negative charges.  Let \(A\) and \(B\) be the largest positive
charge and the largest absolute negative charge used in the relation.
Order its multiset of charges by choosing a negative charge when the
current partial sum is positive, and a positive charge otherwise.  Such a
choice is always available until the sum reaches zero.  Every partial sum
lies in the integer interval \([1-B,A]\).  The partial sums before the
terminal zero are distinct: a repetition would give a proper nonempty
zero-sum subrelation.  There are \(A+B\) integers in that interval, so the
length is at most \(A+B\), which implies \eqref{eq:charge-bound}.

For parity, split every charge atom into its even or odd total \(T\)-degree
class.  Any even-parity sum of charge atoms is a sum of even atoms and
pairs, with repetition, of odd atoms.  These elements generate the
contraction semigroup.  Removing decomposable candidates therefore gives
its full Hilbert basis \(H_S\).

Wick contraction gives the exact moment polynomial
\begin{equation}\label{eq:sparse-moment}
 \E(P_S^m)
 =
 \sum_{\substack{|\nu|=m,\ I(\nu)=J(\nu)\\K(\nu)\ {\rm even}}}
 \frac{m!}{\prod_s\nu_s!}\,
 I(\nu)!\,(K(\nu)-1)!!\prod_sc_s^{\nu_s}.
\end{equation}
Thus the truncated full-support coefficient scheme is most naturally
defined directly on the coefficient torus by the Laurent ideal
\begin{equation}\label{eq:laurent-ideal}
 J_{S,M}
 =
 \left\langle\E(P_S),\ldots,\E(P_S^M)\right\rangle
 \subseteq \mathbb Q[c_s^{\pm1}:s\in S].
\end{equation}
The equations are homogeneous, so one coefficient may be normalized to
one.  There is a second coefficient symmetry.  The
Gaussian-preserving torus
\[
  W\longmapsto tW,\qquad Z\longmapsto t^{-1}Z,\qquad T\longmapsto T
\]
acts on the coefficient of \(s=W^{i_s}Z^{j_s}T^{k_s}\) with charge
\(q_s=i_s-j_s\).  Thus the coefficient torus carries the action
\[
  (\lambda,t)\cdot c_s=\lambda t^{q_s}c_s,
\]
and every moment equation is invariant under \(t\) and homogeneous under
\(\lambda\).

If a support contains coefficients \(c_a,c_b\) of distinct charges, then
over an algebraic closure every orbit meets the slice \(c_a=c_b=1\):
one solves
\[
  t^{q_b-q_a}=c_a/c_b,\qquad
  \lambda=(t^{q_a}c_a)^{-1}.
\]
For \(d=q_b-q_a\), useful invariant Laurent ratios are
\[
  y_s=c_s^d c_a^{q_s-q_b}c_b^{q_a-q_s}
  \qquad(s\neq a,b).
\]
On the slice they become \(y_s=c_s^d\), so these ratios define a finite
map from the slice torus.  The elimination itself uses the free slice
coordinates \(c_s\), rather than assuming that the displayed ratios
generate the full invariant ring.  Consequently the four-term elimination has
only two active coefficient variables whenever two charges occur.  If all
charges agree, the torus acts through overall scaling; such an active
non-one-sided support necessarily has charge zero, and the calculation
falls back to one coefficient normalization.

There are
\[
 \sum_{r=1}^4\binom{35}{r}=59{,}535
\]
raw supports.  We identify a support with its reflection under
\(W\leftrightarrow Z\).  Two elementary filters precede the Groebner
calculation.  If \(H_S\) has one element, then \(\Gamma_S\) has rank one
and the moment at the degree of its generator consists of a single
nonzero coefficient monomial; the coefficient torus contains no solution.
More generally, contraction vectors are enumerated from \(H_S\) through
order twenty; any degree containing exactly one vector gives the same
rejection before factorial Wick weights are expanded.  This cutoff limits
only a sufficient rejection filter.  Supports with
the same ordered charge/parity signature share this enumeration.  Only
the remaining supports are expanded using \eqref{eq:sparse-moment}.  If no nonzero
contraction occurs, all charges \(i_s-j_s\) have one strict sign and the
mixed-moment conclusion follows by bounded weight.  Possible additional
factorizations among disjoint primitive contractions are not used in the
certificate.

For completeness, the latter test is finite in this range.  A zero charge
contracts in one or two copies, according to its \(T\)-parity.  Opposite
nonzero charges \(q>0>r\), with \(|q|,|r|\leq4\), have a primitive
zero-sum relation of length
\[
 \frac{q+|r|}{\gcd(q,|r|)}\leq8;
\]
doubling corrects odd \(T\)-parity.  Hence every support which is not
strictly one-sided has a contraction by order sixteen.

\section{Exact enumeration}

After the preceding filters, the four-term stratum contains \(5{,}718\)
coefficient-torus ideals requiring elimination.  Each is computed in the
Laurent ring of the two-normalization slice when its charge data have rank
two, with the one-normalization fallback just described.  The implementation works in the polynomial ring of the free slice
coordinates and saturates by their product.  This computes the contraction
of the Laurent ideal, so its unit-ideal test is exactly emptiness on the
coefficient torus.  Saturation uses ideal quotients rather than an
auxiliary inverse variable.
\begin{samepage}
Exact Groebner bases over \(\mathbb Q\) for \(J_{S,20}\) leave precisely
the supports
\[
\begin{split}
 &\{Z,ZT^2,WZ^2,W^3T\},\\
 &\{ZT,Z^2,WT,W^2\},\\
 &\{ZT,ZT^3,WZ^2T,W^3\}.
\end{split}
\]
\end{samepage}
The first and third acquire the unit ideal after adjoining moments through
order twenty-four.  The same exact search in the one-, two-, and three-term
strata leaves no active one- or two-term family and two three-term
families.  Including the four-term survivor, the genuine active null
families are
\begin{align}
 P_{\rm lin}&=T+aZ+bW,&1+2ab&=0,                 \label{eq:null-linear}\\
 P_{\rm sh}&=Z+aWT-\frac12a^2W^3,                \label{eq:null-shear}\\
 P_{\rm pr}&=ZT+aZ^2+\frac1{2a^2}WT
                    -\frac1{4a^3}W^2,\quad a\neq0. \label{eq:null-product}
\end{align}
For the three survivors only, clearing Laurent monomials exports triangular
lexicographic certificates.  With the first coefficient normalized to one,
their nontrivial equations are respectively
\[
 1+2c_1c_2,\qquad
 2c_2+c_1^2,\qquad
 2c_1^2c_2-1,\quad c_3+c_1c_2^2.
\]
The finite equations supply necessary coefficient conditions for any
polynomial whose moments vanish at every order.  Conversely, the next
section proves that every member of these three families has the required
all-order behavior.  Reflections \(W\leftrightarrow Z\), nonzero overall
scalings and the coefficient-torus normalizations preserve both pure
vanishing and the mixed-moment conclusion, so no coefficient or support
orbit is lost.  The zero polynomial satisfies the conclusion directly.

\section{All-order interpretation of the survivors}

We first isolate the bounded-weight argument used in all cases.
If every monomial of \(R\) has strictly positive charge \(i-j\), then
every monomial of \(R^m\) has charge at least \(m\).  For a fixed polynomial
\(Q\), let \(B\) be the maximum absolute charge of its monomials, taking
\(B=0\) if \(Q=0\).  For \(m>B\), no monomial of \(QR^m\) has charge zero,
and hence \(\E(QR^m)=0\).  The same proof applies to strictly negative
charges.  In particular it applies to \(R=Z\) and \(R=TZ\).

The Gaussian functional is invariant under complex orthogonal linear
substitutions: the Wick formula expresses every polynomial moment as a
sum of products of the covariance pairing.  This algebraic invariance
does not require the substitution to preserve the real integration
contour.  It allows the following changes of Witt frame.

The family \eqref{eq:null-linear} is a nonzero isotropic linear form.
A complex orthogonal change takes it to a scalar multiple of \(Z\), and
the bounded-weight argument applies.

For \eqref{eq:null-shear}, define
\[
 D=W\partial_T-T\partial_Z,\qquad N=WD.
\]
Then \(D(W)=0\), \(D(T)=W\), and \(D(Z)=-T\).  Thus
\[
 \exp(-aN)(Z)=Z+aWT-\frac12a^2W^3=P_{\rm sh}.
\]
The derivation \(N\) is locally nilpotent: assigning weights \(0,1,2\) to
\(W,T,Z\), respectively, shows that it strictly lowers this nonnegative
weight.  Its exponential \(\phi=\exp(-aN)\) is therefore a polynomial
automorphism.

Gaussian integration by parts in these coordinates gives
\[
 \E(WF)=\E(\partial_ZF),\qquad
 \E(TF)=\E(\partial_TF).
\]
Consequently
\[
 \E(NF)
 =\E(W^2\partial_TF)-\E(WT\partial_ZF)
 =\E(\partial_Z^2\partial_TF)-\E(\partial_T\partial_Z^2F)=0.
\]
Applying this to \(N^{r-1}F\) for \(r\geq1\), and using the finite
exponential expansion on any polynomial, proves
\(\E(\phi(F))=\E(F)\).  For fixed \(Q\),
\[
 \E(QP_{\rm sh}^m)
 =\E\bigl(\phi^{-1}(Q)Z^m\bigr).
\]
The right side vanishes eventually by bounded weight.

Finally,
\begin{equation}\label{eq:product-factorization}
 P_{\rm pr}
 =
 \left(Z+\frac1{2a^2}W\right)
 \left(T+aZ-\frac1{2a}W\right).
\end{equation}
For the covariance pairing \(\langle W,Z\rangle=1\),
\(\langle T,T\rangle=1\), the second factor has self-pairing
\(1+2a(-1/(2a))=0\).  The two factors are orthogonal because
\(-1/(2a)+a/(2a^2)=0\).  The first has nonzero self-pairing \(1/a^2\).
Normalize it to a unit vector; its orthogonal complement is a
nondegenerate two-dimensional space containing the second, isotropic
factor.  Completing a Witt basis in that complement takes the first
factor to \(T\) and the second to a nonzero scalar multiple of \(Z\).
Thus \(P_{\rm pr}\) becomes a scalar multiple of \(TZ\), to which bounded
weight applies.

These arguments also prove the vanishing of every positive pure moment
of the three families, by taking \(Q=1\).  Combined with the exhaustive
finite exclusions and the one-sided supports, they prove
Theorem~\ref{thm:sparse}.

\section{Computation and scope}

The supplied source archive \texttt{source.zip} includes the supplement
\emph{Exact Bounded Sparse Exclusion}, with primary checker
\texttt{verify\_three\_real\_gmc\_sparse\_pareto.py} and helper
\texttt{sparse\_supports.py}.  Their respective SHA256 digests are
\begin{quote}
\footnotesize\ttfamily
5733030a5bee28d0a56d248468f9fbcebe5f76f8844cf0ba2adfc2895097aa9f\\
d573772c0d771a7ffab701a73fa1b0c0c118a36c1b86d2a6e880dc0db114d224
\end{quote}
The checker enumerates supports, applies the charge/parity filters, and
saturates exact rational moment ideals on coefficient slices.  It exports
the three coefficient conditions and verifies
\eqref{eq:product-factorization}; software requirements and replay
instructions accompany the code.  No finite-field output is used to
certify the characteristic-zero result.

The finite exclusions supply necessary conditions for pure vanishing;
Section~4 proves the all-order mixed conclusion.  A replay verifies this
implementation, while the coverage and invariance arguments remain part
of the mathematical proof.  Both bounds---degree at most four and at most
four Witt monomials---apply simultaneously.  Denser cubics and higher-degree
four-term polynomials lie outside the theorem.

\paragraph{AI assistance and review.}
ChatGPT (OpenAI) assisted with mathematical arguments, literature checks,
exposition, proof review and verification code.

\bibliographystyle{plain}
\bibliography{references}

\end{document}
